\section{Abschlussorgane bei Speicherkraftwerken}
\subsection{Schützen}
    \begin{tabular}{@{}ll}
        \textbf{Aufgabe:}                       & Binäre Regelung des Wasserdurchflusses („Auf und zu“)\\
        \textbf{Anwendung:}                     & Grundablass, Seeabschluss, Abschluss gegen das Unterwasser\\
        \textbf{\cgn{Vorteile:}}   & Einfacher Aufbau, sichere Absperrung\\
        \textbf{\crd{Nachteile:}}    & Keine Zwischenstellungen („halb offen“ nicht möglich)\\
        \textbf{Beispiel:}                      & Regulierschütz und Reserveschütz beim Grundablass\\
        \textbf{Sonstiges:}                     & Zwei Schützen pro Grundablass für Betrieb und Revision; \\
                                                & schnelle Schliess- und Öffnungszeiten sind entscheidend, \\
                                                & um Druckstösse zu vermeiden.\\
    \end{tabular}

\subsection{Klappen}
    \begin{tabular}{@{}ll}
        \textbf{Aufgabe:}                       & Wasserfluss absperren (Notverschluss, Wartung)\\
        \textbf{Anwendung:}                     & Saugrohrklappen beim Kraftwerk Mapragg\\
        \textbf{\cgn{Vorteile:}}   & Schnelles Schliessen, einfacher Aufbau\\
        \textbf{\crd{Nachteile:}}    & Keine Regelungsfunktion („halb offen“ nicht möglich)\\
        \textbf{Beispiel:}                      & Saugrohrklappen Mapragg\\
        \textbf{Sonstiges:}                     & Dienen als definitiver Abschluss bei Stillstand oder Wartung; \\
                                                & Schliess- und Öffnungszeiten dürfen nicht verändert werden, \\
                                                & um Druckstösse zu vermeiden.\\
    \end{tabular}

\subsection{Drosselklappen}
    \begin{tabular}{@{}ll}
        \textbf{Aufgabe:}                       & Regelbarer Durchfluss durch horizontales Verschieben \\
                                                & (Spalt lässt Wasser durch)\\
        \textbf{Anwendung:}                     & Einsatz bei variablen Durchflüssen\\
        \textbf{\cgn{Vorteile:}}   & Regelbarkeit des Durchflusses möglich\\
        \textbf{\crd{Nachteile:}}    & Höhere hydraulische Verluste\\
        \textbf{Beispiel:}                      & Horizontal verschiebbarer Rohrschieber\\
        \textbf{Sonstiges:}                     & Wird nur selten bei Speicher- oder Laufwasserkraftwerken eingesetzt, \\
                                                & da dort meist binäre Regelung bevorzugt wird.\\
    \end{tabular}

\subsection{Kugelschieber}
    \begin{tabular}{@{}ll}
        \textbf{Aufbau:}                        & Drehbare Kugel mit Loch\\
        \textbf{Aufgabe:}                       & Absperren des Wasserflusses mit nahezu verlustfreiem Durchfluss im offenen Zustand\\
        \textbf{Anwendung:}                     & Definitiver Abschluss (Betriebsring und Reservering)\\
        \textbf{\cgn{Vorteile:}}   & Im offenen Zustand nahezu verlustfrei\\
        \textbf{\crd{Nachteile:}}    & Wartung und Betrieb der Dichtung (Betriebsring) aufwändig\\
        \textbf{Beispiel:}                      & Kugelschieber in Wasserkraftwerken mit Betriebsring und Reservering\\
        \textbf{Sonstiges:}                     & Betriebsring sorgt für Abdichtung zwischen Gehäuse und Kugel, \\
                                                & Reservering wird nur bei Revision geschlossen und \\
                                                & gegen Wiederöffnen gesichert.\\
    \end{tabular}

\subsection{Ring- und Eckringschieber}
\begin{tabular}{@{}ll}
    \textbf{Aufgabe:}                       & Steuerung des Wasserflusses innerhalb oder ausserhalb des Rohres\\
    \textbf{Anwendung:}                     & Ringschieber: innerh. des Rohres; Eckringschieber: ausserh. des Rohres\\
    \textbf{\cgn{Vorteile:}}   & Flexibler Einbau, verschiedene Steuerungsmöglichkeiten\\
    \textbf{\crd{Nachteile:}}    & Erhöhter Konstruktions- und Wartungsaufwand\\
    \textbf{Beispiel:}                      & Ringschieber und Eckringschieber bei variablen Wasserführungen\\
    \textbf{Sonstiges:}                     & Ermöglichen Zwischenstellungen (feine Regelung), \\
                                            & können jedoch zu höheren hydraulischen Verlusten führen.\\
\end{tabular}